\section{Mass--delay as a scan-based axiom}
\label{app:mass_delay}

\begin{axiom}[Scan-based mass--delay chain]
Internal algebraic obstructions (nonassociative defects) act as phase-impedance centers that require extra scan cycles to resolve locally. The resulting increase in local scan density induces an effective metric stretching/time delay in the readout layer.
\end{axiom}
\noindent
The chain is
\begin{equation}
\begin{aligned}
\text{nonassociative defect}
&\longrightarrow \text{phase-impedance center}
\\
&\longrightarrow \text{extra scan cycles}
\\
&\longrightarrow \text{increased local scan density}
\\
&\longrightarrow \text{metric stretching / time delay (Shapiro-like)}.
\end{aligned}
\end{equation}

The relativistic dispersion $E^2=p^2+m^2$ provides a standard orthogonal decomposition between observable momentum and a rest-energy term. Empirically, gravitational time delay provides an operational benchmark for slowing effects \cite{Shapiro1964}.

\paragraph{External benchmark: Schwarzschild lapse and Shapiro delay.}
In Schwarzschild coordinates, the metric has $g_{tt}=-(1-2GM/(rc^2))c^2$, so a static clock at radius $r$ satisfies
\begin{equation}
\diff\tau=\sqrt{1-\frac{2GM}{rc^2}}\,\diff t,
\end{equation}
exhibiting gravitational time dilation through the lapse factor $\sqrt{1-2GM/(rc^2)}$ \cite{Wald1984}.

Under the interface identification in Proposition~\ref{prop:omega_gr_lapse_interface}, the corresponding routing-cost overhead is
\begin{equation}
\frac{\kappa(r)}{\kappa_0}=\left(1-\frac{2GM}{rc^2}\right)^{-1/2}
=1+\frac{GM}{rc^2}+O\!\left(\frac{G^2M^2}{r^2c^4}\right).
\end{equation}

For light propagation past a gravitating body, the Shapiro time delay for a radar signal is, at leading post-Newtonian order,
\begin{equation}
\Delta t_{\mathrm{Shapiro}}\;\approx\;\frac{2GM}{c^3}\log\!\left(\frac{4r_1 r_2}{b^2}\right),
\end{equation}
for endpoints at radii $r_1,r_2$ and impact parameter $b$ \cite{Shapiro1964,Will1993}.

These standard formulas provide the external operational target for any scan-based identification of ``extra cycles'' with effective delays.
